\section{Background}

Performance estimation is valuable for both ML system developers and infrastructure providers. For system developers, it enables rapid evaluation of the performance of the systems they are building. As model sizes grow, many system parameters require careful tuning, for instance, selecting an appropriate parallelization strategy~\cite{shoeybi2019megatron, zheng2022alpa}. Being able to estimate the performance of different strategies makes it easier to identify the most efficient option. For infrastructure providers, performance estimation allows planning for future hardware deployments.

Due to the accuracy limitations of performance modeling (e.g., roofline \cite{roofline}), developers have turned to static workload simulation for performance estimation. 
Today, two workload simulation methods exist.
The first method is trace-based simulation~\cite{astrasim, astrasim2,daydream,dpro,flexflow,distsim,bang2024vtrain,gui2025multiverse}. A trace is collected in a real execution, and a workload can be extracted from a collected trace. Another method is SimAI~\cite{wang2025simai}, where workloads are generated via implementing a mocked version of the ML frameworks.
Workload simulation leads to accurate performance estimation and is increasingly adopted by industry~\cite{sridharan2023chakra}. 

\begin{figure}[t]
    \centering
    \includegraphics[width=.47\textwidth]{figures/background.pdf}
    \caption{Scheduling logic needs to be reimplemented in current simulators for converting high level workload to detailed computation and communication events.}
    \vspace{-4mm}
    \label{fig:background}
\end{figure}

\parab{Why do static workload simulators need to reimplement ML frameworks?}
For trace-based simulators, there are two fundamental reasons why reimplementation of the ML frameworks is necessary.
First, the key goal of simulation is to explore alternative configurations. Yet, a collected trace reflects only the specific configuration used during execution. To enable configuration exploration, trace-based simulation requires a workload extraction step that lifts the trace into abstract workload, revealing higher-level configurations from actual traces (\autoref{fig:background}). This extraction step typically relies on human understanding of ML framework execution, effectively constructing a reversed version of the framework’s logic. Heuristics can be brittle and may fail to generalize across frameworks; therefore, trace-based simulation often relies on extra manual configurations to help with workload extraction, which increases manual effort.
More importantly, after the user changes the configuration in abstract workload, the simulator needs to turn user's configurations into a detailed execution plan (events) for actual simulation. This scheduling also needs to correctly reflect the framework's scheduling being simulated, which is essentially a reimplementation of the framework's logic.

Beyond trace-based simulation, SimAI \cite{wang2025simai} adopts an alternative approach. Rather than extracting workloads from collected traces, it uses mocked frameworks to directly produce low-level events. However, while this removes the burden of workload extraction, reimplementing scheduling remains necessary in the mocked frameworks. It tightly couples SimAI to specific framework and versions: whenever the underlying ML framework evolves or new framework appears, SimAI’s logic must be updated accordingly to ensure correctness. 

Due to the need of reimplemeting frameworks, both trace-based simulators and SimAI struggle to provide complete support of features in current frameworks and new frameworks. For example, none of the existing simulators support TorchTitan~\cite{liang2025torchtitan}, and none of them can analyze throughput and memory usage of selective activation checkpointing~\cite{korthikanti2023ac,liang2025torchtitan} at the same time. Simulators may even fail to generate the exact same model as the framework. For example, when given the same configuration matching Llama2 7B \cite{touvron2023llama}, the size of model generated by SimAI differs by 7.4\% from native \texttt{GPTModel} in Megatron \cite{shoeybi2019megatron} model library.


\parab{Our approach: Hybrid simulation.} Our objective is straightforward: to build a general simulator in which ML frameworks’ source code can be directly reused. Modern ML frameworks already provide well-maintained implementations and highly configurable parallelization strategies. Leveraging them directly makes the simulator both easier to maintain and more broadly applicable across different frameworks. To achieve this, our approach is to construct a GPU cluster execution environment that closely mirrors a real one, allowing ML frameworks to run on top of the simulated cluster with no modification. Next, we discuss the research challenges to realize this approach and the overview of our design.
